\begin{solucion}

Sea $\alpha = \sqrt{2 + \sqrt{2}}$. Para determinar el grado de la extensión $\mathbb{Q}(\alpha)$, buscaremos su un subcuerpo intermedio que nos permita aplicar el Teorema de la Torre.
Observemos que $\mathbb{Q}(\alpha)$ contiene a $\sqrt{2}$, ya que:
\[
\alpha^2 = 2 + \sqrt{2} \quad \Rightarrow \quad \sqrt{2} = \alpha^2 - 2 \in \mathbb{Q}(\alpha).
\]
Por lo tanto, $\mathbb{Q}(\sqrt{2})$ es un subcampo de $\mathbb{Q}(\alpha)$. El polinomio mínimo de $\sqrt{2}$ sobre $\mathbb{Q}$ es $x^2 - 2$, que es irreducible (ya que $2$ no es un cuadrado racional), así que:
\[
[\mathbb{Q}(\sqrt{2}) : \mathbb{Q}] = 2.
\]
Por lo que si consideramos ahora la extensión $\mathbb{Q}(\sqrt{2}, \alpha)$. Dentro de $\mathbb{Q}(\sqrt{2})$, el elemento $\alpha$ satisface:
\[
\alpha^2 = 2 + \sqrt{2},
\]
es decir, $\alpha$ es raíz del polinomio:
\[
f(x) = x^2 - (2 + \sqrt{2}) \in \mathbb{Q}(\sqrt{2})[x].
\]
Supongamos que $\alpha \in \mathbb{Q}(\sqrt{2})$, es decir, que existen $a, b \in \mathbb{Q}$ tales que $\alpha = a + b\sqrt{2}$. Elevando al cuadrado:
\[
\alpha^2 = a^2 + 2ab\sqrt{2} + 2b^2.
\]
Igualamos con $2 + \sqrt{2}$ y separamos en partes racionales e irracionales:
\[
\begin{cases}
a^2 + 2b^2 = 2,\\
2ab = 1.
\end{cases}
\]
De la segunda ecuación, $a = \frac{1}{2b}$ (con $b \ne 0$). Sustituyendo en la primera:
\[
\left(\frac{1}{2b}\right)^2 + 2b^2 = 2 \quad \Rightarrow \quad \frac{1}{4b^2} + 2b^2 = 2.
\]
Multiplicamos por $4b^2$:
\[
1 + 8b^4 = 8b^2 \quad \Rightarrow \quad 8b^4 - 8b^2 + 1 = 0.
\]
Esta es una ecuación cuadrática en $u = b^2$:
\[
8u^2 - 8u + 1 = 0,
\]
con soluciones:
\[
u = \frac{8 \pm \sqrt{64 - 32}}{16} = \frac{8 \pm \sqrt{32}}{16} = \frac{4 \pm 2\sqrt{2}}{8}.
\]
Este valor no es racional, por lo tanto $b^2 \notin \mathbb{Q}$ y hemos llegado a una contradicción. Así, $\alpha \notin \mathbb{Q}(\sqrt{2})$.

Entonces $f(x)$ es irreducible en $\mathbb{Q}(\sqrt{2})[x]$, y el grado de la extensión es:
\[
[\mathbb{Q}(\sqrt{2}, \alpha) : \mathbb{Q}(\sqrt{2})] = 2.
\]
La cadena de cuerpos es:
\[
\mathbb{Q} \subset \mathbb{Q}(\sqrt{2}) \subset \mathbb{Q}(\alpha).
\]
Por la multiplicación de grados:
\[
[\mathbb{Q}(\alpha) : \mathbb{Q}] = [\mathbb{Q}(\alpha) : \mathbb{Q}(\sqrt{2})] \cdot [\mathbb{Q}(\sqrt{2}) : \mathbb{Q}] = 2 \cdot 2 = 4.
\]


% Partiendo de $\alpha = \sqrt{2 + \sqrt{2}}$:

% \[
% \begin{aligned}
% \alpha^2 &= 2 + \sqrt{2}, \\
% \alpha^2 - 2 &= \sqrt{2}, \\
% (\alpha^2 - 2)^2 &= 2, \\
% \alpha^4 - 4\alpha^2 + 4 &= 2, \\
% \alpha^4 - 4\alpha^2 + 2 &= 0.
% \end{aligned}
% \]

% Entonces $\alpha$ es raíz del polinomio $x^4 - 4x^2 + 2$. Este polinomio es irreducible en $\mathbb{Q}[x]$ (no tiene raíces racionales y pasa el criterio de Eisenstein con $p = 2$ sobre el cambio $x \mapsto x^2$), confirmando que:
% \[
% [\mathbb{Q}(\alpha) : \mathbb{Q}] = 4.
% \]

% \subsection*{Base de la extensión}

% Una base válida para $\mathbb{Q}(\alpha)$ sobre $\mathbb{Q}$ es:
% \[
% \{\,1,\ \sqrt{2},\ \alpha,\ \alpha\sqrt{2}\,\},
% \]
% producto de las bases de $\mathbb{Q}(\sqrt{2})$ sobre $\mathbb{Q}$ y de $\mathbb{Q}(\alpha)$ sobre $\mathbb{Q}(\sqrt{2})$. Alternativamente, se puede usar la base:
% \[
% \{\,1,\ \alpha,\ \alpha^2,\ \alpha^3\,\},
% \]
% derivada del polinomio mínimo de grado $4$.
\end{solucion}